\documentclass[11pt,letterpaper]{article}
\usepackage{technical-design-spec}
\usepackage{longtable,booktabs,array,tabularx}
\title{Secure AWS Platform: Six Architecture Acceptance Gates}
\author{Platform Engineering / Security Architecture}
\date{30 October 2026 --- Proposed}
\begin{document}
\maketitle
\tableofcontents

\section{Purpose and source authority}
This proposed baseline interprets the Reusable Secure AWS Cloud Environment Architecture Guide (28 October 2026) using the Secure AWS Terraform Repository and Deployment Architecture Standard. Mike Cohn's \emph{User Stories Applied} informs backlog quality and task decomposition; it is not a specification of AWS services. No production acceptance or infrastructure deployment is claimed.

\section{Three-repository source of truth}
\begin{description}
\item[latex-docs] Architecture specifications, source-controlled \LaTeX{}, diagrams, governance and verified PDFs.
\item[secure-aws-cloud-platform] Terraform state bootstrap, modules, independent environments, CI controls and infrastructure tests.
\item[taiga-automation] Versioned JSON backlog, story/task reconciliation and GitHub-hosted import workflows.
\end{description}
The documentation pipeline must not assume AWS deploy roles. Infrastructure applies must use independently reviewed GitHub environments, OIDC/STS identities and account allowlists. Taiga creation/import requires a separate authorization.

\section{Current implementation baseline}
The Terraform repository contains a development root, a starter private VPC module, and a non-deploying validation workflow. State backend, deploy roles, staged environment roots, production protections and operational modules remain \emph{planned}. Static validation is not evidence of security control effectiveness.

\section{Architecture invariants}
\begin{enumerate}
\item Separate production account, deployment role, state key and approval trust boundary.
\item Least-privilege workload identities; MFA/federation for humans and short-lived OIDC for CI.
\item Explicit network egress and ingress policy; private subnet location is not sufficient.
\item Encryption, state lock and protected restore evidence before privileged changes.
\item Typed Terraform module contracts, reviewed changes, and negative security tests.
\item No plaintext credentials in Git, ordinary environment files, published PDFs or CI logs.
\end{enumerate}

\section{Acceptance gate evidence}
\begin{longtable}{@{}p{0.08\linewidth}p{0.30\linewidth}p{0.53\linewidth}@{}}
\toprule Gate & Outcome & Minimum objective evidence \\ \midrule \endhead
G1 & Foundation and isolation & Account and role boundaries reviewed; negative cross-environment permission tests recorded \\
G2 & Networking and IAM & Intended-only flows and least-privilege policies checked with negative access tests \\
G3 & State and secrets & Versioned and encrypted remote state; locking and secret isolation verified \\
G4 & Module reuse & Same reviewed Terraform module revision composes isolated dev/staging roots \\
G5 & Monitoring and recovery & Alerts assigned to owners and restore drill achieves reviewed recovery objectives \\
G6 & CI/CD governance & PR static checks, protected delivery, OIDC roles and independent production approval verified \\
\bottomrule\end{longtable}
\section{Decision and exclusions}
Every gate requires a human review of linked implementation tests, negative security assertions and risk exceptions. Start with a single Lambda or ECS/Fargate workload selected by ADR. EKS, service meshes, elaborate transit routing and cross-Region active-active recovery remain out of initial scope.
\end{document}